\section{Introduction}\label{sec:introduction}

Superposing two independent dimer coverings produces loops whose large-scale
geometry is expected to be conformally invariant. For the square lattice, the
predicted limit is the conformal loop ensemble with parameter \(4\). The
distinction between geometric and topological convergence is important here:
knowing which punctures a loop surrounds does not control thin excursions,
invisible retraced pieces, or repeated traversal of a limiting trace.
We prove the curve convergence for the standard Temperleyan law in the upper
half-plane, retaining every nested generation.

\subsection{The law and the topology}
Write
\[
 \HH=\{z\in\mathbb C:\operatorname{Im}z>0\},\qquad
 \HH_\delta=\HH\cap\delta\mathbb Z^2 .
\]
Edges join nearest neighbors. A dimer covering is a perfect matching.
To specify its infinite-volume law, start with a simple lattice polygon in
the coarse grid \((0,\delta)+2\delta\mathbb Z^2\), retain the fine-grid
vertices and edges in its closure, and delete one coarse boundary vertex,
called the root. Uniform matchings on these finite Temperleyan graphs
converge locally as the polygons exhaust \(\HH\) and their roots escape to
infinity. We use this standard law, equivalently specified by the inverse
Kasteleyn operator vanishing at infinity \cite[Section~2]{BI}.
The coarse-grid translate puts the lower boundary on the bottom row of
\(\HH_\delta\).

Take two independent matchings with this law. Almost surely their
superposition consists of doubled edges and finite simple even cycles
\cite[opening of Section~2]{BI}; see also \cite[Section~5.4, Lemma~26]{Dubedat}.
Delete the doubled edges and let
\(\mathcal L_\delta\) be the collection of all cycles, drawn with straight
edges and regarded as unrooted, unoriented loops.
Nested \(\CLE_4\) means an ordinary nonnested \(\CLE_4\), followed recursively
by conditionally independent copies inside every loop. The normalization is
the one in which the outermost ensemble comes from a Brownian loop soup of
central charge \(1\) \cite{SW}.

Let \(\DD=\{w\in\mathbb C:|w|<1\}\). We fix, throughout, the compactification
\[
 F:\HH\longrightarrow\DD,\qquad F(z)=\frac{z-i}{z+i}.
\]
For continuous loops \(\gamma,\widetilde\gamma:S^1\to\overline{\DD}\), put
\[
 d_{\rm loop}(\gamma,\widetilde\gamma)
 =\inf_{\phi}\sup_{t\in S^1}
       |\gamma(t)-\widetilde\gamma(\phi(t))|,
\]
where \(\phi\) ranges over circle homeomorphisms of either orientation.
We identify loops at zero distance; in particular a weakly monotone change
of parameter, with pauses, does not change the loop.
Collections omit constant loops and have finitely many members of diameter
greater than any fixed positive number, with multiplicities retained.
An \(\eps\)-matching of two
collections is a partial bijection covering every loop of diameter greater
than \(\eps\) on either side, with matched loops at distance at most
\(\eps\). The unmatched loops therefore all have diameter at most
\(\eps\). Convergence means that such matchings exist with
\(\eps\downarrow0\). All diameters and loop distances in this definition
are measured after applying \(F\).

It is convenient to use a compatible metric \(\rho\) on collections.
For a partial bijection, charge the distance of each matched pair and
half the diameter of each unmatched loop; take the supremum of these
charges, then the infimum over partial bijections.
This defines \(\rho\), with the supremum of an empty family equal to zero.
The triangle inequality follows by composing partial bijections: a loop
whose partner becomes unmatched has half-diameter at most the sum of
the two charges, since
\[
 |\diam\gamma-\diam\widetilde\gamma|
       \le 2d_{\rm loop}(\gamma,\widetilde\gamma).
\]
Vanishing \(\rho\) identifies equal collections because there are only
finitely many loops above each positive diameter. A matching of cost at
most \(\eps\) leaves only loops of diameter at most \(2\eps\) unmatched,
so \(\rho\) induces exactly the topology specified above.
We use the inherited topology on the subspace of collections whose loops
are contained in the open disk. This is the target space for the theorem.
The larger closed-disk space will be useful during the compactness proof,
before boundary contact has been excluded.

\begin{theorem}\label{thm:main}
For the half-plane Temperleyan law specified above, the family
\(F(\mathcal L_\delta)\) is tight as \(\delta\downarrow0\) in the
loop-collection topology just defined, and
\[
 F(\mathcal L_\delta)\ \Longrightarrow\ F(\mathcal L),
 \qquad \delta\downarrow0,
\]
where \(\mathcal L\) is standard nested \(\CLE_4\) in \(\HH\).
The convergence holds along the full mesh limit and matches all
macroscopic loops as curves.
\end{theorem}

\subsection{History and the remaining geometric question}
Theorem~\ref{thm:main} resolves the half-plane Temperleyan form of the
double-dimer \(\CLE_4\) scaling-limit conjecture. Two strands of earlier
work explain both its prediction and the remaining obstacle.

Exact dimer enumeration on the square lattice was developed by
Kasteleyn~\cite{Kasteleyn} and Temperley and Fisher~\cite{TemperleyFisher},
using Pfaffian and determinant methods. Lieb~\cite{Lieb} introduced the
transfer-matrix approach. These finite algebraic descriptions provide
the background for our rectangle calculation. At the continuum scale,
Kenyon~\cite{KenyonGFF} proved convergence of appropriately normalized
dimer-height fluctuations for suitable Temperleyan approximations of
smooth Jordan domains to the Gaussian free
field. That field limit motivates the connection with \(\CLE_4\), but
does not by itself give convergence of loop curves: passing from a
distribution-valued field to its interfaces requires additional control.

Further evidence for the interface prediction came from Kenyon and
Wilson~\cite{KW}. For alternating black and white
boundary nodes in the half-plane, they showed that the limiting
double-dimer pairing probabilities agree with those for the corresponding
Gaussian-free-field contour lines. This identifies the
boundary pairings, while leaving convergence of the paths open.
Kenyon~\cite{Kenyon} used quaternionic weights to construct loop-homotopy
observables with conformally invariant scaling limits. Dub\'edat~\cite{Dubedat} identified
a class of these topological correlators with their nested \(\CLE_4\)
counterparts through isomonodromic tau-functions.
Basok and Chelkak~\cite[Corollary~1.7]{BC} recovered probabilities of cylindrical events
from their lamination expansion, subject to a continuum lamination-tail
estimate; Bai and Wan~\cite[Corollary~1.5]{BW} supplied that estimate
for the nested ensemble.
We use the precise half-plane formulation in Basok and Izyurov's
preprint~\cite[Theorem~4.2(2), Eq.~(4.17)]{BI}.
It retains all nested generations. Basok and Izyurov also prove local
Gaussian fluctuations and convergence of the centered nesting field,
which describe further aspects of the ensemble beyond a fixed puncture
record.

The distinction needed here is between those puncture records and
convergence in the curve topology. A long thin excursion may surround
none of the selected points. Moreover, even equality of a limiting trace
and winding number one allows a path to move backward along an arc before
continuing forward. We must control both effects while retaining the
number of loops and their nesting. The new estimates address these
geometric questions directly.

\subsection{The geometric estimates and the proof}

The new input consists of two geometric estimates in a corner-rooted
Temperleyan rectangle. The first bounds the total number of disjoint
traversals of a fixed-width slab, including multiple traversals by a single
loop. It supplies both macroscopic loop counts and moduli of continuity
after reparametrization. The second controls the signed crossing counts
of arcs whose endpoints reach prescribed opposite sides of the slab;
the sign records the direction of a crossing. On a shrinking transverse
interval it excludes an arc's nonzero count disagreeing with that of its
whole loop, and simultaneous nonzero counts from arcs on distinct loops.
After localization, this prevents cancellation of an excursion against a
return path that finite-puncture data would fail to see.

Both estimates come from a transfer matrix counting dimer configurations
column by column. A boundary state records which dimers cross the seam
between columns. Complementing alternate occupancy bits turns such a
state into a subset of transverse sites, called a particle state; the
column transfer is an exterior power of a one-particle matrix.
This uses Lieb's transfer method~\cite{Lieb}, with sine modes and
particle sectors corresponding to the open-boundary spectral analysis of
Rasmussen and Ruelle~\cite[Section~IV~A]{RR}. We derive the needed formulas in
the convention appropriate to the deleted corner and prove the cap
normalization explicitly, where a cap is the part of the rectangle on
one side of a chosen seam.
Changing particle number by a large amount gives a quadratic spectral
penalty and hence traversal tightness. A two-particle insertion supported
on a short interval moves two occupied sites from one matching's state to
the other's. Its matrix coefficients are two-by-two minors
of an interval projection. Modes spread across the transverse sites make these coefficients
quadratic in the interval length. To turn that signed insertion into a
probability bound, we prove a pointwise positivity identity, weighting
nested loop portions cut off by transverse lines, so that contributions
decrease geometrically down each
nesting chain. This separation of a local operator estimate from a planar
positivity argument is the main technical mechanism. The estimates are
proved in Sections~\ref{sec:transfer} and~\ref{sec:short-segment}.

Wilson's algorithm and the Temperley correspondence~\cite{Wilson,PW,KPW}
transfer the estimates
to bounded half-plane windows, uniformly in the mesh
(Section~\ref{sec:localization}). That section then proves compactness by
constructing parametrizations directly from the fixed-slab estimates.
The passage from crossing counts to controlled parametrizations follows
the method of Aizenman and Burchard~\cite[Lemma~2.4 and Section~4]{AB}.

Section~\ref{sec:identification} couples a subsequential path limit and
its puncture records with a canonical nested \(\CLE_4\). It fixes each
finite puncture grid before choosing a sufficiently small lattice mesh.
On the resulting joint realization, a pathwise argument rules out
invisible curves and identifies each visible trace. A second grid
refinement then excludes backtracking along that trace. This last step
recovers the traversal order that trace equality and winding number one
leave undetermined.

\subsection{Established inputs}\label{sec:inputs}
We state precisely the information imported from the literature. A
lamination relative to a finite set of punctures records the homotopy
classes of disjoint simple loops in the punctured domain, including their
multiplicities. Delete every loop surrounding at most one puncture.
The remaining record will be called the reduced lamination.

\begin{theorem}[Cylindrical convergence {\cite[Theorem~4.2(2)]{BI}}]
\label{thm:cylindrical}
Let \(z_1,\ldots,z_r\) be distinct interior points of \(\HH\), and choose
face punctures \(z_j^\delta\to z_j\). For the two independent matchings
used in Theorem~\ref{thm:main}, the reduced lamination relative to
\((z_j^\delta)_{j=1}^r\) converges in law to the corresponding reduced
lamination of nested \(\CLE_4\). In particular the multisets of enclosed
puncture subsets, with multiplicities and with subsets of size at most
one deleted, converge in law.
\end{theorem}

The precise input is equation~(4.17) of the cited theorem in the
January~2, 2025 preprint version. It is uniform when the puncture tuples remain in a compact
set away from collisions. We only need the weaker formulation above.
For a fixed puncture set the lamination state space is countable.
Convergence of its point probabilities to a probability law implies total
variation convergence: first restrict to a finite set carrying nearly all
limiting mass. Passing to enclosed subsets is then a measurable
pushforward; we do not claim that subsets determine every homotopy class.
Later refinements always fix a finite puncture configuration before taking
the lattice mesh sufficiently small.

\begin{proposition}[Standard properties of nested \(\CLE_4\)]
\label{prop:cle-facts}
After conformal transport to the disk, nested \(\CLE_4\) consists almost
surely of pairwise disjoint simple loops contained in the open disk.
There are finitely many loops of diameter greater than any positive
number. Every fixed interior point is almost surely off all loop traces
and is surrounded by infinitely many successive nested loops.
\end{proposition}

The nonnested properties and the recursive construction are established in
\cite[Sections~2.1--2.2]{SW}. In particular, every fixed interior point is
surrounded almost surely; conditional iteration in the successive loop
interiors gives infinitely many surrounding loops and avoidance of all
traces. Children lie strictly inside their parents, so disjointness also
persists across generations. Bounded-domain macroscopic finiteness for the
entire nested ensemble is recorded explicitly in \cite[p.~2788]{BC}.
We also use the classical Temperley correspondence \cite[Theorem~1]{KPW}
and Wilson's rooted spanning-tree algorithm \cite{Wilson}, in the
order-independent form of \cite[Theorems~13 and~16]{PW}.
Their exact application to the finite rectangles will be stated when the
localization argument is introduced.

\paragraph{Conventions.}
Transfer lengths and transverse site counts use lattice units; geometric
windows and buffers use physical coordinates. Thus the transverse size is
an odd integer \(n=2p+1\), and a transfer through \(t\) columns has physical
width \(\delta t\). Fix the checkerboard bipartition into black and white
vertices. A cycle is directed by following first-matching edges from
black to white and second-matching edges from white to black. All signed
intersection sums use a common transverse orientation. Their sign changes
when the direction of the cycle is reversed, while every event we use is
invariant under that reversal. Constants may depend on fixed rectangles,
windows, and positive buffers, but not on the vanishing mesh.
